\documentclass[conference,onecolumn,12pt]{IEEEtran}

\usepackage{setspace}
\onehalfspacing
\usepackage[T1]{fontenc}
\usepackage[utf8]{inputenc}
\usepackage{amsmath}
\usepackage{amssymb}
\usepackage{booktabs}
\usepackage{array}
\usepackage{cite}
\usepackage{url}
\usepackage{graphicx}
\usepackage{hyperref}
\hypersetup{hidelinks}

\begin{document}

% ─── Title ───────────────────────────────────────────────────────────────────

\title{Traffic Congestion Detection and Classification with\\
Sensor-Noise Robustness: A Digital Twin-Enabled\\
Multi-Model Machine Learning Approach}

% ─── Authors ─────────────────────────────────────────────────────────────────

\author{
  \IEEEauthorblockN{%
    Assem Mohamed Saad$^{1}$,\;
    Omar Ahmed Asaad$^{1}$,\;
    Habiba Hosam Eldin$^{1}$\\[2pt]
    Seif Adel Elbasha$^{1}$,\;
    Owida Elsayed$^{1}$
  }
  \IEEEauthorblockA{%
    \{200030013, 200027710, 200025297, 200046355, 200042768\}@must.edu.eg
  }
  \IEEEauthorblockN{%
    Dr.\ Amr Ibrahim$^{1}$,\;
    Dr.\ Walid Hamdy$^{1}$
  }
  \IEEEauthorblockA{%
    \{amr.ibrahim,\;walid.hamdy\}@must.edu.eg
  }
  \IEEEauthorblockA{%
    $^{1}$Faculty of Information Technology,
    Misr University for Science and Technology (MUST), Giza, Egypt
  }
}

\maketitle

% ─── Abstract ────────────────────────────────────────────────────────────────

\begin{abstract}
Urban traffic congestion is one of the most persistent challenges facing
modern smart cities, imposing significant economic, environmental, and social
costs. Despite growing adoption of data-driven approaches for traffic
monitoring, most existing systems treat sensor data as noise-free and
reliable---an assumption that fails in real-world deployments where cameras
suffer from occlusion and sensors experience dropout events. This paper
presents a traffic congestion detection and classification system that
combines multiple machine learning models with a digital twin-inspired
simulation layer to evaluate model robustness under realistic sensor
degradation conditions. Using a two-month dataset of vehicle counts sampled
at 15-minute intervals, we train and compare five classifiers---Logistic
Regression, Random Forest, Extra Trees, Histogram Gradient Boosting, and a
fully-connected Neural Network---on a four-class congestion labeling task
(normal, low, high, heavy). Beyond standard clean-data evaluation, we
systematically inject simulated occlusion noise ($\pm$3 vehicle count error)
and sensor dropout (2\% of readings) to benchmark each model's resilience.
Temporal features are encoded cyclically to preserve time-of-day continuity,
and a strict chronological train/test split is enforced to prevent data
leakage. A robustness decay metric---the macro-averaged F1 drop from clean to
noisy conditions---is reported alongside accuracy, balanced accuracy,
precision, and recall for all five models. All trained artifacts, confusion
matrices, and robustness decay visualizations are saved for full
reproducibility.
\end{abstract}

\begin{IEEEkeywords}
traffic congestion classification, sensor noise robustness, digital twin
simulation, machine learning, ensemble methods, neural networks, intelligent
transportation systems
\end{IEEEkeywords}

% ─── 1. Introduction ─────────────────────────────────────────────────────────

\section{Introduction}

The rapid growth of urban populations has intensified the demand for
intelligent transportation systems (ITS) capable of monitoring, classifying,
and responding to traffic conditions in real time. Traffic congestion not only
wastes fuel and increases travel times, but also contributes
disproportionately to carbon emissions and urban air
pollution~\cite{abdallah2026}. As cities invest in sensor
infrastructure---loop detectors, camera-based vehicle counters, acoustic
sensors, and GPS probes---the volume of available traffic data has grown
substantially, creating an opportunity for machine learning-based congestion
classification systems.

However, translating this data richness into reliable real-world systems
remains non-trivial. Sensor deployments are inherently imperfect: cameras are
occluded by large vehicles or adverse weather, sensor connections drop
intermittently, and count readings are subject to systematic
biases~\cite{xie2024,liu2024}. The majority of published ML systems for
traffic classification evaluate only on clean test sets, ignoring the
performance degradation that occurs once deployed in the field. This gap
between benchmark accuracy and operational reliability motivates a more
rigorous evaluation protocol.

Parallel to this, the concept of the digital twin---a virtual replica of a
physical system that mirrors its real-time state---has gained traction as a
framework for testing and validating decisions before they are applied to real
infrastructure~\cite{bagabaldo2025,kaewunruen2022}. By simulating sensor
failure scenarios within a digital twin environment, we can stress-test
classification models without waiting for failures to occur in production.

This paper addresses both challenges. We build a multi-model traffic
congestion classification system trained on two months of real sensor data and
evaluate each model under both clean and synthetically degraded conditions.
The degradation simulation mirrors known failure modes: random count
perturbations (occlusion) and full zero-readings (sensor dropout). We compare
five classifiers spanning traditional, ensemble, and neural network families,
and report a robustness decay metric---the F1-Macro drop from clean to noisy
conditions---as a primary evaluation criterion alongside standard accuracy
metrics.

The remainder of this paper is organized as follows.
Section~\ref{sec:problem} formally states the problem.
Section~\ref{sec:contributions} describes our contributions relative to
existing work. Section~\ref{sec:related} reviews related literature.
Section~\ref{sec:methodology} details our proposed methodology.
Section~\ref{sec:results} presents results and discussion.
Section~\ref{sec:conclusion} concludes the paper and outlines future
directions.

% ─── 2. Problem Statement ────────────────────────────────────────────────────

\section{Problem Statement}
\label{sec:problem}

Let the feature vector at time step $t$ be
\[
  \mathbf{x}_t = [\text{CC},\;\text{BC},\;\text{BuC},\;\text{TC},\;
                  \text{HS},\;\text{HC},\;\mathbf{d}],
\]
where CC, BC, BuC, and TC are vehicle counts (car, bike, bus, and truck),
HS and HC are cyclic hour projections, and $\mathbf{d}$ is a one-hot
encoding of the day of the week. The
classification target
$y_t \in \{\text{normal},\,\text{low},\,\text{high},\,\text{heavy}\}$
represents the congestion level observed at that time step.

The dataset $\mathcal{D} = \{(\mathbf{x}_t, y_t)\}_{t=1}^{T}$ consists of
observations sampled at 15-minute intervals over a two-month period, yielding
$T \approx 5{,}952$ records and 96 samples per day. The class distribution is
imbalanced, with \emph{normal} comprising approximately 61\% of records and
\emph{high} only 6\%.

Two evaluation regimes are defined:
\begin{itemize}
  \item \textbf{Clean regime:} Models are evaluated on the held-out test set
    $\mathcal{D}_{\text{test}}$ with original sensor readings.
  \item \textbf{Noisy regime:} A perturbed test set
    $\tilde{\mathcal{D}}_{\text{test}}$ is constructed by applying additive
    uniform noise $\epsilon \sim \mathcal{U}(-3,\,3)$ to each vehicle count
    and randomly zeroing 2\% of count readings to simulate dropout.
\end{itemize}

The core challenge is to learn a classifier $f\colon \mathbf{x} \mapsto y$
that achieves high F1-Macro on $\mathcal{D}_{\text{test}}$ while minimizing
the \emph{robustness decay}:
\begin{equation}
  \Delta F1 \;=\; F1^{\text{clean}}_{\text{macro}} - F1^{\text{noisy}}_{\text{macro}}
  \label{eq:decay}
\end{equation}

A system that scores well on clean data but degrades sharply under noise is
unsuitable for real-world deployment. A secondary challenge is the time-series
nature of the data: random splitting would leak future information into
training. We therefore perform a chronological split at the day boundary,
reserving the last 20\% of days for testing.

% ─── 3. Contributions ────────────────────────────────────────────────────────

\section{Contributions}
\label{sec:contributions}

Relative to the existing literature reviewed in Section~\ref{sec:related},
this work contributes the following:

\textbf{1.~Robustness-first evaluation protocol.}
Most published traffic classification systems report accuracy on clean,
i.i.d.\ test sets~\cite{kashyap2023,mahmoud2024,sharma2024}. We introduce a
paired evaluation---clean vs.\ synthetically noisy---that directly quantifies
each model's resilience to sensor degradation. This protocol is inspired by
digital twin simulation principles~\cite{abdallah2026,kaewunruen2022} and is
absent from the majority of ML-for-traffic benchmarks.

\textbf{2.~Simulation of physically grounded sensor failure modes.}
Rather than applying generic Gaussian noise, our perturbation model
specifically targets occlusion ($\pm$3 count error per vehicle type) and
sensor dropout (2\% zero-out rate)---failure modes documented in real sensor
deployments~\cite{xie2024,liu2024}. This grounding distinguishes our
simulation from abstract noise injection schemes.

\textbf{3.~Multi-family model comparison on a real-world tabular traffic
dataset.}
We compare five classifiers from three algorithmic families (linear, tree
ensemble, neural network) on identical feature sets and splits. While general
tabular benchmarks exist~\cite{mcelfresh2023}, comparative traffic-specific
benchmarks using balanced accuracy and macro-averaged metrics alongside
robustness measures are sparse~\cite{ting2021,sharma2024}.

\textbf{4.~Temporally faithful feature engineering and train/test splitting.}
We encode time-of-day as sine/cosine projections to preserve cyclic continuity
and enforce a strict chronological split at the day level---practices not
consistently adopted in traffic ML literature~\cite{kashyap2023,chen2025}.

\textbf{5.~Reproducible artifact pipeline.}
All trained models, scalers, label encoders, confusion matrices, and
robustness decay plots are saved programmatically, enabling direct
reproduction and deployment.

% ─── 4. Related Work ─────────────────────────────────────────────────────────

\section{Related Work}
\label{sec:related}

\subsection{Traffic Congestion Detection and Classification}

Traffic flow modeling has evolved from classical parametric methods toward
data-driven approaches. A comprehensive survey by Kashyap et
al.~\cite{kashyap2023} categorizes the field into \emph{prediction}
(estimating future flow values) and \emph{classification} (assigning discrete
congestion labels) tasks, noting that classification is better suited to
real-time signal control and alert systems. Our work falls squarely in the
classification category, targeting a four-level taxonomy also adopted in recent
IEEE and Springer studies~\cite{mahmoud2024,sharma2024}.

Inputs to traffic classification systems vary considerably across the
literature. Image-based systems use CNN and YOLO architectures to detect
vehicles directly from camera feeds~\cite{chen2025}, while sensor-based
systems operate on pre-aggregated count or speed data~\cite{mahmoud2024}.
Mahmoud et al.~\cite{mahmoud2024} train classifiers on sensor-derived vehicle
counts---the same input modality as our system---and report that Random Forest
consistently outperforms SVM and Logistic Regression baselines. A comparative
study~\cite{sharma2024} confirms that ensemble methods are competitive with
deep sequence models (LSTM, GRU) on structured tabular traffic data,
particularly when training sets are limited in size. Our dataset uses vehicle
counts across four categories (car, bike, bus, truck) at 15-minute
granularity, closely mirroring the input structure of IoV-based monitoring
systems~\cite{ali2021}. Class imbalance is addressed through class-weighted
objectives and balanced accuracy reporting.

\subsection{ML Model Comparison: Ensembles vs.\ Neural Networks}

A long-standing open question in applied ML is when neural networks outperform
gradient-boosted decision trees on tabular data. McElfresh et
al.~\cite{mcelfresh2023} benchmark 19 algorithms across 176 datasets at
NeurIPS 2023, finding that the NN vs.\ GBDT performance gap is frequently
negligible and that light GBDT hyperparameter tuning often matters more than
model choice---directly motivating our head-to-head comparison on traffic
data.

Ting et al.~\cite{ting2021} compare Random Forest against multiple graph
convolutional network (GCN) variants for traffic prediction, finding that RF
remains a strong baseline even against spatially-aware deep models. In our
setting, which lacks graph structure (no road network topology), tree
ensembles are expected to be especially competitive. We extend the comparison
to Extra Trees, Histogram Gradient Boosting, and a regularized
fully-connected neural network to provide a broad view of the performance
landscape on this specific task.

\subsection{Digital Twin Simulation for Traffic Management}

The digital twin paradigm---maintaining a synchronized virtual replica of a
physical system for monitoring, prediction, and control---has seen growing
adoption in urban traffic
management~\cite{abdallah2026,bagabaldo2025,kaewunruen2022,alrakhami2025}.
Abdallah and Alghamdi~\cite{abdallah2026} demonstrate a GRU-based digital
twin for real-time traffic signal optimization, achieving a 33\% improvement
in predictive accuracy and a 15\% reduction in CO$_2$ emissions. Their
architecture couples IoT sensor feeds with a simulation model that validates
control actions before applying them to the physical network---a design
principle we adopt at the evaluation level by stress-testing classifiers
within a simulated degraded sensor environment.

Bagabaldo and Hackl~\cite{bagabaldo2025} review over 70 papers on digital
twins for intelligent intersections, identifying three key benefits: testing
interventions without real-world risk, generating synthetic training data, and
detecting anomalies by comparing predicted vs.\ observed system states. Our
noise injection module operationalizes the second benefit. Kaewunruen and
Lian~\cite{kaewunruen2022} describe the real-time synergy between live data
streams and simulation models as the core digital twin architectural pattern,
while Al-Rakhami and Al-Mashari~\cite{alrakhami2025} identify traffic
management as the highest-priority application domain for urban digital twin
deployments.

\subsection{Sensor Noise and Robustness in Traffic Monitoring}

Real-world traffic sensor deployments are subject to multiple degradation
mechanisms. Xie et al.~\cite{xie2024} document that distributed acoustic
sensing (DAS) systems are susceptible to polarization fading and quasi-random
optical noise, with accuracy dropping substantially under low
signal-to-noise ratio conditions---motivating systematic robustness
evaluation across all sensor modalities. Liu et al.~\cite{liu2024} show that
dropout regularization and noisy data augmentation improve generalization under
deployment-time distribution shift, providing the closest methodological
precedent to our noise injection protocol.

At the general ML level, Ibias et al.~\cite{ibias2024} develop a theoretical
framework for noise robustness through input abstractions, demonstrating
significantly less performance degradation under additive and dropout noise
for models with abstraction layers. While their experiments focus on network
traffic classification, the analytical framework applies directly to our
sensor-count inputs, and their findings suggest that ensemble methods' implicit
averaging confers natural robustness---consistent with our experimental design.

% ─── 5. Proposed Methodology ─────────────────────────────────────────────────

\section{Proposed Methodology}
\label{sec:methodology}

Our system is composed of four sequential stages: data preprocessing and
feature engineering, chronological train/test splitting with feature scaling,
sensor noise simulation, and multi-model training with paired evaluation.

\subsection{Dataset}

The dataset is a real-world traffic sensor log collected over two months at a
monitored road segment, recorded at 15-minute intervals (96 samples per day,
$\approx$5,952 total records after cleaning). Each record contains:
\emph{Time}, \emph{Date}, \emph{Day of the week}, \emph{CarCount},
\emph{BikeCount}, \emph{BusCount}, \emph{TruckCount}, \emph{Total}, and the
target label \emph{Traffic Situation}
$\in \{\text{normal},\,\text{low},\,\text{high},\,\text{heavy}\}$.

The class distribution is imbalanced: \emph{normal} accounts for 3,610
records ($\approx$61\%), \emph{heavy} for 1,137 ($\approx$19\%), \emph{low}
for 834 ($\approx$14\%), and \emph{high} for only 371 ($\approx$6\%).

\subsection{Data Preprocessing and Feature Engineering}

Raw data is deduplicated and rows with missing values are dropped. The
\texttt{Time} column is parsed in 12-hour format to extract the hour of day
as an integer in $[0,\,23]$. To capture the cyclic nature of time---where
hour 23 is adjacent to hour 0---the hour is projected onto the unit circle:
\begin{align}
  \text{Hour\_Sin} &= \sin\!\left(\frac{2\pi\cdot\text{Hour}}{24}\right), \notag\\
  \text{Hour\_Cos} &= \cos\!\left(\frac{2\pi\cdot\text{Hour}}{24}\right)
  \label{eq:cyclic}
\end{align}

This two-dimensional representation ensures that the distance between adjacent
hours is uniform across the midnight boundary, which a raw integer encoding
would distort~\cite{kashyap2023}. The \texttt{Day of the week} column is
one-hot encoded into seven binary indicator variables (all seven days retained
to avoid ambiguity in the multi-class context). The final feature vector is:
\begin{equation}
  \mathbf{x} =
  [\,\text{HS},\;\text{HC},\;\text{CC},\;\text{BC},\;\text{BuC},\;\text{TC},\;
   \mathbf{d}_{\text{Mon}},\;\ldots,\;\mathbf{d}_{\text{Sun}}\,]
  \label{eq:features}
\end{equation}
where HS = Hour\_Sin, HC = Hour\_Cos, CC = CarCount,
BC = BikeCount, BuC = BusCount, TC = TruckCount. The \texttt{Total}
column is excluded from the feature set to avoid a direct linear proxy for the
label. The target is integer-encoded via a \texttt{LabelEncoder}, mapping the
four string classes to indices 0--3.

\subsection{Chronological Train/Test Split and Feature Scaling}

To respect the temporal ordering of the data and prevent any form of data
leakage, the dataset is split chronologically at day boundaries. The total
number of complete days is $n_{\text{days}} = \lfloor T / 96 \rfloor$, and
the split point is set such that the last 20\% of days form the test set:
\begin{equation}
  \text{split\_day}
    = n_{\text{days}} - \max\!\bigl(1,\;\mathrm{round}(0.2\times n_{\text{days}})\bigr)
  \label{eq:split}
\end{equation}

All rows before \texttt{split\_day} constitute the training set; all remaining
rows form the test set. Random shuffling is explicitly avoided. Feature
standardization is applied using a \texttt{StandardScaler} fitted exclusively
on the training set and then applied to both splits---preventing test-set
statistics from influencing the scaling~\cite{mcelfresh2023}.

\subsection{Sensor Noise Simulation}

To emulate the degradation conditions encountered in real-world deployments, a
noisy test set $\tilde{\mathcal{D}}_{\text{test}}$ is constructed using two
independent perturbations applied via a seeded random number generator
(seed\,=\,42 for reproducibility):

\textbf{Occlusion noise.} Each of the four vehicle count features
(CarCount, BikeCount, BusCount, TruckCount) is perturbed by an additive
integer drawn uniformly from $\{-3,\ldots,+3\}$, then clipped at zero to
prevent physically impossible negative values. This simulates the systematic
undercounting that occurs when vehicles partially occlude one another in a
camera's field of view~\cite{xie2024,liu2024}.

\textbf{Sensor dropout.} For 2\% of randomly selected records, one vehicle
count column chosen uniformly at random is set to zero, simulating complete
sensor loss for a single vehicle category---a failure mode observed in both
loop detector malfunctions and camera feed interruptions~\cite{xie2024}.

Noise is applied \emph{only} to the test set. The training set remains clean,
reflecting the realistic scenario where models are trained on curated
historical data but deployed against imperfect live feeds. Both
$\mathcal{D}_{\text{test}}$ and $\tilde{\mathcal{D}}_{\text{test}}$ share
identical labels $y_{\text{test}}$; only the feature values differ.

\subsection{Model Architectures and Hyperparameters}

Five classifiers spanning three algorithmic families are trained on the same
standardized training set and evaluated on both test conditions:

\textbf{Logistic Regression (LR).} SAGA solver, 2,000 maximum iterations,
balanced class weights (random state\,=\,42). Serves as a linear lower-bound
reference.

\textbf{Random Forest (RF).} 300 decision trees grown with bootstrap
sampling, balanced class weights, full parallelization
(random state\,=\,42). Averaging over many trees provides implicit variance
reduction expected to translate into noise robustness~\cite{ting2021}.

\textbf{Extra Trees (ET).} 600 extremely randomized trees, balanced class
weights (random state\,=\,42). Unlike RF, split thresholds are selected
randomly rather than optimally, further reducing variance; the larger
ensemble size (600 vs.\ 300 in RF) compensates for higher individual-tree
error.

\textbf{Histogram Gradient Boosting (HGB).} Maximum tree depth of 6
(random state\,=\,42). Discretizes continuous features into histograms for
efficient training. Class imbalance is handled implicitly through the boosting
loss landscape, as this implementation does not support a \texttt{class\_weight}
parameter directly.

\textbf{Neural Network (NN).} A fully-connected feedforward network with the
following architecture:
\begin{align*}
  &\text{Input}
   \to \mathrm{Dense}(128,\text{ReLU})
   \to \mathrm{BN}
   \to \mathrm{Dropout}(0.3) \\
  &\quad\to \mathrm{Dense}(64,\text{ReLU})
   \to \mathrm{BN}
   \to \mathrm{Dense}(32,\text{ReLU}) \\
  &\quad\to \mathrm{Dense}(K,\text{Softmax})
\end{align*}
where $K = 4$ is the number of congestion classes. Compiled with the Adam
optimizer and sparse categorical cross-entropy loss. Training runs for up to
100 epochs with a batch size of 32 and a 10\% validation split. Early stopping
monitors validation loss with a patience of 10 epochs and restores the best
weights upon termination. The Dropout and Batch Normalization layers improve
generalization under the noise-induced distribution shift between clean and
noisy test sets~\cite{ibias2024}.

\subsection{Evaluation Protocol and Metrics}

Each of the five models is evaluated on both $\mathcal{D}_{\text{test}}$
(clean) and $\tilde{\mathcal{D}}_{\text{test}}$ (noisy), producing two
prediction vectors per model. The following metrics are computed for each
(model, setting) pair:

\begin{itemize}
  \item \textbf{Accuracy:} fraction of correctly classified samples.
  \item \textbf{Balanced Accuracy (BalAcc):} mean per-class recall,
    compensating for class imbalance.
  \item \textbf{F1-Macro:} unweighted mean of per-class F1 scores---the
    primary metric for imbalanced classification~\cite{kashyap2023}.
  \item \textbf{Precision-Macro / Recall-Macro:} unweighted per-class means.
  \item \textbf{Robustness Decay ($\Delta F1$):} as defined in
    Equation~\eqref{eq:decay}; lower is better.
\end{itemize}

Confusion matrices are generated for all 10 (model $\times$ setting)
combinations and saved as heatmap visualizations using the \texttt{magma}
colormap. A horizontal bar chart ranks all five models by $\Delta F1$ to
provide an at-a-glance robustness comparison.

\subsection{Artifact Saving}

All outputs are written to \texttt{outputs/traffic\_final\_research/} for
reproducibility. Table~\ref{tab:artifacts} lists every saved file.

\begin{table}[h]
  \centering
  \caption{Saved Artifacts and Filenames}
  \label{tab:artifacts}
  \renewcommand{\arraystretch}{1.2}
  \begin{tabular}{ll}
    \toprule
    \textbf{Artifact} & \textbf{Filename} \\
    \midrule
    Trained Keras neural network  & \texttt{traffic\_nn\_model.keras} \\
    Fitted StandardScaler         & \texttt{traffic\_scaler.pkl} \\
    Fitted LabelEncoder           & \texttt{traffic\_encoder.pkl} \\
    Confusion matrices (10 total) & \texttt{cm\_\{model\}\_\{setting\}.png} \\
    Robustness decay chart        & \texttt{robustness\_decay.png} \\
    Full metrics table            & \texttt{full\_experiment\_results.csv} \\
    \bottomrule
  \end{tabular}
\end{table}

% ─── 6. Results and Discussion ───────────────────────────────────────────────

\section{Results and Discussion}
\label{sec:results}

\textit{This section will be completed after experiments are executed.
Placeholder tables are provided below to define the reporting structure.}

\subsection{Classification Performance---Clean Setting}

Table~\ref{tab:clean} will report all five metrics for each model evaluated on
the clean test set $\mathcal{D}_{\text{test}}$.

\begin{table}[h]
  \centering
  \caption{Classification Performance on Clean Test Data}
  \label{tab:clean}
  \renewcommand{\arraystretch}{1.2}
  \begin{tabular}{lccccc}
    \toprule
    \textbf{Model} & \textbf{Acc.} & \textbf{BalAcc} & \textbf{F1} & \textbf{Prec.} & \textbf{Rec.} \\
    \midrule
    Logistic Reg.  & --- & --- & --- & --- & --- \\
    Random Forest  & --- & --- & --- & --- & --- \\
    Extra Trees    & --- & --- & --- & --- & --- \\
    Hist.\ GBM     & --- & --- & --- & --- & --- \\
    Neural Network & --- & --- & --- & --- & --- \\
    \bottomrule
  \end{tabular}
\end{table}

\subsection{Classification Performance---Noisy Setting}

Table~\ref{tab:noisy} will report the same five metrics for all models
evaluated on the noisy test set $\tilde{\mathcal{D}}_{\text{test}}$, which
applies occlusion ($\pm$3 count perturbation) and sensor dropout (2\%
zero-out rate).

\begin{table}[h]
  \centering
  \caption{Classification Performance on Noisy Test Data}
  \label{tab:noisy}
  \renewcommand{\arraystretch}{1.2}
  \begin{tabular}{lccccc}
    \toprule
    \textbf{Model} & \textbf{Acc.} & \textbf{BalAcc} & \textbf{F1} & \textbf{Prec.} & \textbf{Rec.} \\
    \midrule
    Logistic Reg.  & --- & --- & --- & --- & --- \\
    Random Forest  & --- & --- & --- & --- & --- \\
    Extra Trees    & --- & --- & --- & --- & --- \\
    Hist.\ GBM     & --- & --- & --- & --- & --- \\
    Neural Network & --- & --- & --- & --- & --- \\
    \bottomrule
  \end{tabular}
\end{table}

\subsection{Robustness Decay Analysis}

Table~\ref{tab:decay} will summarize the robustness decay ($\Delta F1$,
Equation~\eqref{eq:decay}) for each model, ranked from most to least robust.
The accompanying bar chart generated by the pipeline will provide a visual
ranking.

\begin{table}[h]
  \centering
  \caption{Robustness Decay---F1-Macro Drop (Clean $\to$ Noisy)}
  \label{tab:decay}
  \renewcommand{\arraystretch}{1.2}
  \begin{tabular}{lccc}
    \toprule
    \textbf{Model} & \textbf{F1 (Clean)} & \textbf{F1 (Noisy)} & $\boldsymbol{\Delta F1}$ $\downarrow$ \\
    \midrule
    Logistic Reg.  & --- & --- & --- \\
    Random Forest  & --- & --- & --- \\
    Extra Trees    & --- & --- & --- \\
    Hist.\ GBM     & --- & --- & --- \\
    Neural Network & --- & --- & --- \\
    \bottomrule
  \end{tabular}
\end{table}

\subsection{Per-Class Analysis and Confusion Matrices}

Confusion matrices for all 10 (model $\times$ setting) combinations will be
analyzed here, with particular attention to the minority classes \emph{high}
and \emph{low}, which are most vulnerable to noise-induced misclassification.
The effect of class-weighted training (LR, RF, ET) on minority-class recall
will be contrasted against the unweighted HGB model.

\subsection{Discussion}

This subsection will interpret the results in light of the related work
reviewed in Section~\ref{sec:related}. Specifically, it will address:
(i) whether the NN vs.\ GBDT performance gap reported by McElfresh et
al.~\cite{mcelfresh2023} holds under this traffic classification task;
(ii) whether ensemble averaging translates to lower robustness decay as
suggested by Ibias et al.~\cite{ibias2024}; and (iii) how our system
compares to the GRU-based digital twin approach of Abdallah and
Alghamdi~\cite{abdallah2026} in terms of classification accuracy and
operational reliability under sensor degradation.

% ─── 7. Conclusion ───────────────────────────────────────────────────────────

\section{Conclusion and Future Work}
\label{sec:conclusion}

This paper presented a traffic congestion classification system grounded in
real sensor data, evaluated under both clean and simulated noisy conditions to
bridge the gap between benchmark performance and operational robustness. Five
classifiers across three algorithm families were trained and compared on a
four-class congestion labeling task, with robustness decay as a primary metric
alongside standard accuracy measures. A digital twin-inspired noise injection
layer simulates physically grounded sensor failure modes---camera occlusion and
dropout---enabling a rigorous pre-deployment stress test absent from most
existing traffic ML benchmarks.

Future work will extend this system in three directions: (1) integrating
road-network topology to enable spatially-aware modeling using graph neural
networks; (2) replacing static noise injection with a learned digital twin
model that adapts its failure simulation to observed degradation statistics;
and (3) deploying the best-performing classifier in a real-time edge inference
system with continuous retraining as new sensor data arrives.

% ─── References ──────────────────────────────────────────────────────────────

\begin{thebibliography}{00}

\bibitem{abdallah2026}
W. Abdallah and M. Alghamdi,
``Digital twin-enabled AI for sustainable traffic management: real-time urban
mobility optimization in smart cities,''
\emph{PeerJ Comput.\ Sci.}, vol.~12, p.~e3574, 2026.
doi:~\href{https://doi.org/10.7717/peerj-cs.3574}{10.7717/peerj-cs.3574}

\bibitem{ali2021}
F.~Ali, A.~Ali, M.~Imran, R.~A.~Naqvi, M.~H.~Siddiqi, and K.-S.~Kwak,
``Machine learning approach on traffic congestion monitoring system in Internet
of Vehicles,''
\emph{J.\ Inf.\ Process.\ Syst.}, vol.~17, no.~6, pp.~1207--1224, 2021.
doi:~\href{https://doi.org/10.1016/j.procs.2020.09.212}{10.1016/j.procs.2020.09.212}

\bibitem{alrakhami2025}
M.~S.~Al-Rakhami and M.~Al-Mashari,
``Digital twin technology in smart cities: A step toward intelligent urban
management,''
\emph{Energy Rep.}, vol.~13, pp.~1--15, 2025.
doi:~\href{https://doi.org/10.1016/j.egyr.2025.01.012}{10.1016/j.egyr.2025.01.012}

\bibitem{bagabaldo2025}
A.~R.~Bagabaldo and J.~Hackl,
``Digital twins for intelligent intersections: A literature review,''
\emph{arXiv preprint arXiv:2510.05374}, 2025.
doi:~\href{https://doi.org/10.48550/arXiv.2510.05374}{10.48550/arXiv.2510.05374}

\bibitem{chen2025}
L.~Chen, Y.~Zhang, and X.~Wang,
``Traffic congestion recognition based on convolutional neural networks in
different scenarios,''
\emph{Eng.\ Appl.\ Artif.\ Intell.}, vol.~145, p.~109782, 2025.
doi:~\href{https://doi.org/10.1016/j.engappai.2025.109782}{10.1016/j.engappai.2025.109782}

\bibitem{ibias2024}
A.~Ibias, K.~Capa\l{}a, V.~Varma, A.~Dr\'o\.zd\.z, and J.~L.~R.~Sousa,
``Improving noise robustness through abstractions and its impact on machine
learning,''
\emph{arXiv preprint arXiv:2406.08428}, 2024.
doi:~\href{https://doi.org/10.48550/arXiv.2406.08428}{10.48550/arXiv.2406.08428}

\bibitem{kaewunruen2022}
S.~Kaewunruen and Q.~Lian,
``A digital twin in transportation: Real-time synergy of traffic data streams
and simulation for virtualizing motorway dynamics,''
\emph{Adv.\ Eng.\ Inform.}, vol.~55, p.~101858, 2022.
doi:~\href{https://doi.org/10.1016/j.aei.2022.101858}{10.1016/j.aei.2022.101858}

\bibitem{kashyap2023}
A.~A.~Kashyap, R.~Patil, T.~J.~Rama Sushma, A.~Bhandari, R.~Gupta, and
A.~Bhardwaj,
``A survey on traffic flow prediction and classification,''
\emph{Intell.\ Syst.\ Appl.}, vol.~18, p.~200205, 2023.
doi:~\href{https://doi.org/10.1016/j.iswa.2023.200205}{10.1016/j.iswa.2023.200205}

\bibitem{liu2024}
X.~Liu, R.~Chen, and J.~Park,
``Enhancing traffic monitoring with noise-robust distributed acoustic sensing
and deep learning,''
\emph{Transp.\ Res.\ C, Emerg.\ Technol.}, vol.~165, p.~104756, 2024.
doi:~\href{https://doi.org/10.1016/j.trc.2024.104756}{10.1016/j.trc.2024.104756}

\bibitem{mahmoud2024}
M.~Mahmoud, A.~Alkhatib, and H.~Hamdoun,
``Traffic congestion prediction using machine learning,''
in \emph{Proc.\ IEEE Int.\ Conf.\ Comput.\ Inf.\ (ICCI)}, 2024.
doi:~\href{https://doi.org/10.1109/ICCI.2024.10498418}{10.1109/ICCI.2024.10498418}

\bibitem{mcelfresh2023}
D.~C.~McElfresh, S.~Khandagale, J.~Valverde, C.~VishakPrasad, B.~Feuer,
C.~Hegde, G.~Ramakrishnan, M.~Goldblum, and C.~White,
``When do neural nets outperform boosted trees on tabular data?''
in \emph{Adv.\ Neural Inf.\ Process.\ Syst.\ (NeurIPS)}, 2023.
doi:~\href{https://doi.org/10.48550/arXiv.2305.02997}{10.48550/arXiv.2305.02997}

\bibitem{sharma2024}
R.~Sharma, M.~Gupta, and K.~Patel,
``Intelligent traffic flow prediction using deep learning techniques:
A comparative study,''
\emph{SN Comput.\ Sci.}, vol.~5, no.~4, p.~512, 2024.
doi:~\href{https://doi.org/10.1007/s42979-024-03552-3}{10.1007/s42979-024-03552-3}

\bibitem{ting2021}
T.~J.~Ting, X.~Li, S.~Sanner, and B.~Abdulhai,
``Revisiting random forests in a comparative evaluation of graph convolutional
neural network variants for traffic prediction,''
in \emph{Proc.\ IEEE Int.\ Conf.\ Intell.\ Transp.\ Syst.\ (ITSC)}, 2021,
pp.~1--8.
doi:~\href{https://doi.org/10.1109/ITSC48978.2021.9564595}{10.1109/ITSC48978.2021.9564595}

\bibitem{xie2024}
D.~Xie, X.~Wu, Z.~Guo, H.~Hong, B.~Wang, and Y.~Rong,
``Intelligent traffic monitoring with distributed acoustic sensing,''
\emph{Seismol.\ Res.\ Lett.}, 2024.
doi:~\href{https://doi.org/10.1785/0220240298}{10.1785/0220240298}

\bibitem{zhang2025}
Y.~Zhang, H.~Liu, Z.~Wang, and J.~Chen,
``AI-powered urban transportation digital twin: Methods and applications,''
\emph{arXiv preprint arXiv:2501.10396}, 2025.
\url{https://arxiv.org/abs/2501.10396}

\end{thebibliography}

\end{document}
